% reproducibility.tex -- Appendix B
\label{sec:reproducibility}

This appendix is a compact reproducibility quickstart.  Goal: a
reader who has just cloned the repository at the v1.0 tag can
reproduce every audit-table row of Paper~II and rebuild this
PDF from a fresh checkout in under ten minutes.

\subsection{Environment}
\label{subsec:repro_environment}

\begin{itemize}
\item \textbf{Operating system.}  Verified on Windows~11 (MSYS2
  UCRT64) and any POSIX system with the equivalent toolchain.
  No platform-specific code.
\item \textbf{Python.}  Python~3.10 or later.  The framework's
  symbolic verification scripts use only \texttt{sympy} (no
  numerical-only dependencies); the dependency manifest is
  \nolinkurl{virtual-env-requirements.txt} in the repository
  root.
\item \textbf{LaTeX.}  Any modern TeX distribution (TeX Live~2023+
  or MiKTeX) with \texttt{pdflatex} and standard packages
  (\texttt{amsmath}, \texttt{hyperref}, \texttt{longtable},
  \texttt{caption}, \texttt{adjustbox}; see
  \nolinkurl{paper/macros/packages.tex}).
\item \textbf{GNU Make}~$\geq$ 4.3 (the stock Windows GNU~Make port
  is 3.81 and too old; use the \texttt{chocolatey/make} build or
  MSYS2's \texttt{make}).
\end{itemize}

\subsection{From fresh clone to verified paper}
\label{subsec:repro_steps}

\begin{verbatim}
git clone https://github.com/JackDMenendez/dcl-paper-02-sm-derivation
cd dcl-paper-02-sm-derivation
git checkout v1.0

python -m venv .venv-ucrt64
source .venv-ucrt64/bin/activate    # POSIX / MSYS2
# .venv-ucrt64\Scripts\activate     # Windows native cmd
pip install -r virtual-env-requirements.txt

make paper                          # rebuilds build/Paper.pdf
\end{verbatim}

The \texttt{make paper} target runs three passes of
\texttt{pdflatex} (plus a \texttt{bibtex} pass for the
bibliography), producing
\nolinkurl{build/Paper.pdf} byte-equivalent to the released
artifact archived in
\nolinkurl{.stage/Paper_v1.0.pdf}.

\subsection{Verifying every audit-table row}
\label{subsec:repro_audit}

The end-to-end verification of Paper~II's PASS / PART / FAIL
status rows is running the thirteen scripts in
\nolinkurl{src/utilities/}.  Each script is self-contained and
runs as a Python module:

{\scriptsize
\begin{verbatim}
python -m src.utilities.automorphism_discrete
python -m src.utilities.automorphism_rgb_su3
python -m src.utilities.automorphism_direct_product
python -m src.utilities.automorphism_direct_product_extended
python -m src.utilities.tick_rule_extended_consistency
python -m src.utilities.tick_rule_gauge_invariance
python -m src.utilities.su3_branch_consistency
python -m src.utilities.chirality_parity_alignment
python -m src.utilities.tick_rule_cp_modified
python -m src.utilities.su3_generation_from_colour_memory
python -m src.utilities.automorphism_centralizer_extended
python -m src.utilities.aut_centralizer_extras_commutators
python -m src.utilities.induced_gauge_action_nonabelian
\end{verbatim}
}

All but \nolinkurl{aut_centralizer_extras_commutators.py} complete
in seconds on a consumer-grade laptop.  The bracket-structure
script (1711 commutators of $12\times 12$ sympy matrices) takes
several minutes; this is the most expensive computation in the
paper and is the load-bearing verification for the structural
finding that the 71-dim discrete centralizer is the standard
subalgebra-with-complement of a simple semisimple Lie algebra
(see~\autoref{subsec:terminal_characterisation}).

\subsection{Verifying the central numerical claim}
\label{subsec:repro_central}

The framework's first quantitative gauge-coupling prediction
($g_3^2 / g_2^2 = 3/2$ at the lattice scale, Eq.~\eqref{eq:g3_over_g2})
is derived in \nolinkurl{src/utilities/induced_gauge_action_nonabelian.py}.
The relevant lines extract the prediction from the Q-tensor and
the spectator-factor counting:

{\scriptsize
\begin{verbatim}
import sympy as sp
from sympy import Rational
# Spectator-factor multiplicities on the per-site C^{12}:
N_f_SU2 = 2 * 3   # chirality x colour spectators = 6
N_f_SU3 = 2 * 2   # chirality x isospin spectators = 4
T_F = Rational(1, 2)  # fundamental of SU(N)
inv_g2_sq = N_f_SU2 * T_F   # = 3
inv_g3_sq = N_f_SU3 * T_F   # = 2
print(inv_g2_sq / inv_g3_sq)   # 3/2 (= g_3^2 / g_2^2 at lattice scale)
\end{verbatim}
}

The full script also verifies Paper~I's Q-tensor with eigenvalues
$\{4, 4, 16\}$, the $T_F$ trace normalisation, and the abelian
($U(1)$) case where all 12 components of the per-site amplitude
carry unit charge, giving $1/g_1^2 : 1/g_2^2 : 1/g_3^2 = 12 : 3 : 2$
at the lattice scale.

\subsection{Where to look next}
\label{subsec:repro_pointers}

\begin{itemize}
\item For the per-claim verification trail: navigate from any
  audit-table row in Table~\ref{tab:audit} to the script named in
  its evidence column.  Run that script.
\item For structural-interpretation context: the corresponding
  note in \nolinkurl{notes/} contains the working derivation, the
  representation-theoretic decomposition, and pointers to adjacent
  notes (the upstream-flow tags at the bottom of each note name
  the related findings).
\item For the framework's foundational substrate:
  Paper~I~\cite{menendez2026geometry}.  Section~5--6 of Paper~I
  derive the Dirac structure from the bipartite tick rule;
  Section~7 derives gravity-as-clock-density; Appendix~B
  establishes the bipartite-plaquette Q-tensor and the
  structural form of the induced gauge action.  Paper~II's
  Phase~4 calculation extends Paper~I's Appendix~B unchanged to
  the non-abelian case.
\item For follow-on questions and open items: the catalogue at
  \nolinkurl{external/dcl/notes/follow_on_implications.md}
  (Paper~I repository) enumerates the planned downstream papers;
  Paper~II is item~\#16 of that catalogue.
\end{itemize}

\subsection{Release flow (for reference)}
\label{subsec:repro_release_flow}

The release procedure for v1.0 (and subsequent versions) is
documented in \nolinkurl{release_notes/README.md} of the
repository.  Briefly: (i)~final commits land on \texttt{main};
(ii)~the change log
(\nolinkurl{release_notes/v1.0.md}) and the GitHub Release body
(\nolinkurl{release_notes/v1.0-release-message.md}) are drafted;
(iii)~deposit on Zenodo to receive the DOI; (iv)~commit the
version-bump filling in the DOI in \nolinkurl{paper/main.tex},
\nolinkurl{CITATION.cff}, and the release-notes placeholders;
(v)~rebuild the PDF and snapshot to \nolinkurl{.stage/};
(vi)~tag and push; (vii)~create the GitHub Release.  This appendix
and the \nolinkurl{release_notes/README.md} parent describe the
full reproducibility chain from a fresh clone of any tagged
release.
